\subsection{FIELDS OF POINT DISTRIBUTIONS DEFINED BY THE ACTION OF A GROUP ON A MANIFOLD}

Let \((g,x)\mapsto\lambda(g,x)=gx\) be a law of left operation of class \(C^{r}\) of G on X. Let \(s\leqslant r\) and \(t\in U_{s}(G)\) . For all \(x\in X\) , \(t*\varepsilon_{x}\in T_{x}^{(s)}(X)\) . The mapping \(x\mapsto t*\varepsilon_{x}\) is called the field of point distributions defined by t and the action of G on X and denoted sometimes by \(D_{t}^{\lambda}\) or simply \(D_{t}\) . Let \(\Omega\) be an open subset of X and F a Hausdorff polynormed space. If \(f\colon\Omega\to F\) is of class \(C^{r}\) and \(s\leqslant r\) , the function \(t^{\vee}\ast f\) on \(\Omega\) is also denoted by \(D_{t}.f\) . Then

\[
(D _ {t} f) (x) = \langle t * \varepsilon_ {x}, f \rangle .\tag{16}
\]

If \(s < \infty\) , then \(D_t f \in \mathcal{C}^{r-s}(\Omega, F)\) by no. 4. Thus \(f \mapsto D_t f\) is a mapping of \(\mathcal{C}^r(\Omega, F)\) into \(\mathcal{C}^{r-s}(\Omega, F)\) (often denoted by \(D_t\) by an abuse of notation).

If \(t \in \mathrm{U}_s(\mathrm{G})\) , \(t' \in \mathrm{U}_{s'}(\mathrm{G})\) and \(s + s' \leqslant r\) , then, by Proposition 18 of no. 4,

\[
\mathrm{D} _ {t * t ^ {\prime}} f = \mathrm{D} _ {t ^ {\prime}} (\mathrm{D} _ {t} f)\tag{17}
\]

and hence, using the abuse of notation indicated above,

\[
\mathrm{D} _ {t * t ^ {\prime}} = \mathrm{D} _ {t ^ {\prime}} \circ \mathrm{D} _ {t}.\tag{18}
\]

Suppose that G and X are finite-dimensional. The mapping \((t, x) \mapsto t \otimes \varepsilon_{x}\) of \(\mathrm{T}^{(\mathrm{s})}(\mathrm{G}) \times \mathrm{X}\) into the vector bundle \(\mathrm{T}^{(\mathrm{s})}(\mathrm{G} \times \mathrm{X})\) (cf.~Differentiable and Analytic Manifolds, R, 13.2.5) is of class \(C^{r-s}\) . Hence (Differentiable and Analytic Manifolds, R, 13.2.5) the mapping \((t, x) \mapsto t * \varepsilon_{x}\) of \(\mathrm{T}^{(\mathrm{s})}(\mathrm{G}) \times \mathrm{X}\) into the vector bundle \(\mathrm{T}^{(\mathrm{s})}(\mathrm{X})\) is of class \(C^{r-s}\) . In particular, \(D_{t}\) is a differential operator of order \(\leqslant s\) and class \(C^{r-s}\) in the sense of Differentiable and Analytic Manifolds, R, 14.1.6. By formula (16), the function \(D_{t}f\) is then the image of f under this differential operator (Differentiable and Analytic Manifolds, R, 14.1.4).

We now no longer suppose that G and X are finite-dimensional. Let \(\psi\) be an automorphism of the manifold X and \(\Delta\) a field of point distributions on X. Conforming with the general definitions, the transform of \(\Delta\) under \(\psi\) is the field of point distributions on X whose value at \(\psi(x)\) is \(\psi_{*}(\Delta(x))\) ; we denote this mapping by \(\psi(\Delta)\) . If \(g \in G\) and \(\tau(g)\) denotes the automorphism \(x \mapsto gx\) of X, the transform of \(\Delta\) under \(\tau(g)\) is also called the transform of \(\Delta\) under g.

PROPOSITION 21. Let \(\psi\) be an automorphism of X commuting with the operations of G. Then \(D_{t}\) is invariant under \(\psi\) .

For all \(x \in X\) ,

\[
\begin{array}{r l} (\psi (\mathrm{D} _ {t})) (\psi (x)) & = \psi_ {*} (\mathrm{D} _ {t} (x)) = \psi_ {*} (t * \varepsilon_ {x}) \\ & = t * \psi_ {*} (\varepsilon_ {x}) \quad \text {(Proposition 15)} \\ & = t * \varepsilon_ {\psi (x)} = \mathrm{D} _ {t} (\psi (x)). \end{array}
\]

PROPOSITION 22. If \(g \in G\) , the transform of \(D_t\) under \(g\) is \(D_{\varepsilon_g * t * \varepsilon_{g^{-1}}}\) .

The value of this transform at gx is

\[
\begin{array}{r l r} \tau (g) _ {*} (\mathrm{D} _ {t} (x)) & = \tau (g) _ {*} (t * \varepsilon_ {x}) \\ & = \varepsilon_ {g} * (t * \varepsilon_ {x}) & \text {(Proposition 14 (i))} \\ & = (\varepsilon_ {g} * t * \varepsilon_ {g^{-1}}) * \varepsilon_ {g x} & \text {(Propositions 1 and 2)} \\ & = \mathrm{D} _ {\varepsilon_ {g} * t * \varepsilon_ {g^{-1}}} (g x). \end{array}
\]

Let \((x,g)\mapsto \mu (x,g) = xg\) be a law of right operation of class \(\mathbf{C}^r\) of G on X. Let \(s\leqslant r\) and \(t\in \mathrm{U}_{s}(\mathrm{G})\) . For all \(x\in \mathbf{X}\) , \(\varepsilon_{x}*t\in \mathrm{T}_{x}^{(s)}(\mathbf{X})\) . The mapping \(x\mapsto \varepsilon_x*t\) is called the field of distributions defined by \(t\) and the action of G on X and is sometimes denoted by \(\mathrm{D}_{t}^{\mu}\) or simply \(\mathrm{D}_t\) . Let \(\Omega\) be an open subset of X. If \(f:\Omega \to F\) is of class \(\mathbf{C}^r\) , the function \(f*t^{\vee}\) is denoted by \(\mathrm{D}_t f\) . Then

\[
\left(\mathrm{D} _ {t} f\right) (x) = \langle \varepsilon_ {x} * t, f \rangle\tag{19}
\]

and, in the obvious notation,

\begin{enumerate}
\def\labelenumi{(\arabic{enumi})}
\setcounter{enumi}{19}
\tightlist
\item
\end{enumerate}

\[
D _ {t * t ^ {\prime}} f = \mathrm{D} _ {t} (\mathrm{D} _ {t ^ {\prime}} f)\tag{20}
\]

\[
\mathrm{D} _ {t * t ^ {\prime}} = \mathrm{D} _ {t} \circ \mathrm{D} _ {t ^ {\prime}}\tag{21}.
\]

Proposition 21 remains valid. Let \(g \in G\) . The transform of \(D_t\) under \(g\) (that is under the automorphism \(x \mapsto xg\) of \(X\) ) is \(D_{\varepsilon_{g^{-1}} * t * \varepsilon_g}\) .

